% ============================================================
% Kapitel 4: Evaluation
% ============================================================
\chapter{Evaluation}
\label{sec:evaluation}

Dieses Kapitel evaluiert die im IZPL eingesetzte On-Premise-Sprachpipeline in zwei Bereichen:
die \textbf{LLM-Inferenzleistung} des Qwen 3.5 9B Modells (via Ollama auf einer NVIDIA RTX 3080 Ti)
sowie die \textbf{Spracherkennungsqualität} des Parakeet TDT 0.6B ASR-Modells.
Alle Benchmarks wurden auf dem in Abschnitt~\ref{sec:izpl} beschriebenen Laboraufbau durchgeführt.

% ============================================================
\section{LLM-Benchmark: Qwen 3.5 9B}
\label{sec:eval_llm}
% ============================================================

Zur Evaluierung der LLM-Inferenzleistung wurden fünf robotikspezifische Prompts
aus dem Bereich ROS2 und dem igus ReBeL 6-DOF-Roboter verwendet. Jeder Prompt
wurde einzeln verarbeitet; dabei wurden Token-Generierungsgeschwindigkeit,
Latenz, GPU-/CPU-Auslastung und Energieverbrauch gemessen.

\subsection{Generierungsgeschwindigkeit}

\begin{figure}[H]
    \centering
    \includegraphics[width=0.85\textwidth]{llm_generierungsgeschwindigkeit.png}
    \caption{Token-Generierungsgeschwindigkeit (Tokens/s) für fünf robotikspezifische Prompts.
    Alle Prompts werden mit einer konsistenten Geschwindigkeit von ca.\,93 Tokens/s verarbeitet.}
    \label{fig:eval_gen_speed}
\end{figure}

Die Generierungsgeschwindigkeit liegt konsistent bei durchschnittlich
\textbf{93,7 Tokens/s} über alle fünf Prompts (Abbildung~\ref{fig:eval_gen_speed}).
Die geringe Varianz ($\pm$\,1,5 Tokens/s) zeigt, dass die Inferenzgeschwindigkeit
weitgehend unabhängig von der Prompt-Komplexität ist, sie hängt
primär von der Speicherbandbreite der GPU ab.

\subsection{Antwortzeit-Aufschlüsselung}

\begin{figure}[H]
    \centering
    \includegraphics[width=0.85\textwidth]{llm_antwortzeit.png}
    \caption{Zeitliche Aufschlüsselung der Antwortzeiten: Modell laden, Prompt-Auswertung (Prefill) und Generierung. Die Gesamtzeit skaliert linear mit der Antwortlänge.}
    \label{fig:eval_latency}
\end{figure}

Abbildung~\ref{fig:eval_latency} schlüsselt die Antwortzeit auf.
Die Prefill-Phase (Prompt-Auswertung) benötigt nur wenige Sekunden,
während die Generierungsphase proportional zur Antwortlänge skaliert.
Für den längsten Prompt (ROS2 Node Code, $\approx$10.000 Tokens)
beträgt die Gesamtantwortzeit ca.\,110 Sekunden.

\subsection{GPU-Auslastung und VRAM}

\begin{figure}[H]
    \centering
    \includegraphics[width=0.85\textwidth]{llm_gpu_vram.png}
    \caption{GPU-Auslastung und VRAM-Belegung während des Benchmarks. Die GPU-Auslastung liegt
    stabil bei ca.\,90\,\%, der VRAM-Bedarf beträgt konstant ca.\,8\,GB.}
    \label{fig:eval_gpu_vram}
\end{figure}

Die GPU-Auslastung liegt während der Inferenz stabil bei ca.\,90\,\%
(Abbildung~\ref{fig:eval_gpu_vram}). Der VRAM-Bedarf des Qwen 3.5 9B
beträgt konstant ca.\,8\,GB und liegt damit deutlich unter dem
12-GB-Limit der RTX 3080 Ti. Der verbleibende VRAM steht für den
KV-Cache und ggf. parallele Prozesse zur Verfügung.

\subsection{CPU- und RAM-Auslastung}

\begin{figure}[H]
    \centering
    \includegraphics[width=0.85\textwidth]{llm_cpu_ram.png}
    \caption{CPU- und RAM-Auslastung während der LLM-Inferenz. Die CPU-Belastung bleibt
    unter 15\,\%, der RAM-Bedarf bei ca.\,4\,GB.}
    \label{fig:eval_cpu_ram}
\end{figure}

Da die gesamte Inferenz auf der GPU stattfindet, bleibt die CPU-Auslastung
gering ($<$\,15\,\%, Abbildung~\ref{fig:eval_cpu_ram}). Der Arbeitsspeicherverbrauch
ist mit ca.\,4\,GB ebenfalls moderat, was den parallelen Betrieb weiterer
Anwendungen (z.\,B. ROS2-Nodes) ermöglicht.

\subsection{GPU-Leistungsaufnahme}

\begin{figure}[H]
    \centering
    \includegraphics[width=0.85\textwidth]{llm_gpu_leistung.png}
    \caption{GPU-Leistungsaufnahme während des Benchmarks. Die Leistung liegt bei ca.\,320\,W
    (TDP-Grenze: 350\,W), was einer nahezu vollständigen Auslastung der GPU entspricht.}
    \label{fig:eval_gpu_power}
\end{figure}

Die GPU-Leistungsaufnahme beträgt während der Inferenz ca.\,320\,W
(Abbildung~\ref{fig:eval_gpu_power}), was nahe an der TDP-Grenze von
350\,W liegt. Für den Dauerbetrieb ist daher eine ausreichende Kühlung
sicherzustellen.

\subsection{Netzwerk- und Festplatten-E/A}

\begin{figure}[H]
    \centering
    \includegraphics[width=0.85\textwidth]{llm_netzwerk_io.png}
    \caption{Kumulativer Netzwerk- und Festplatten-E/A während des Benchmarks.
    Der Netzwerkverkehr resultiert aus der Kommunikation mit der Ollama-API.}
    \label{fig:eval_io}
\end{figure}

Abbildung~\ref{fig:eval_io} zeigt den kumulativen E/A-Verlauf.
Da das Modell nach dem initialen Laden vollständig im GPU-Speicher
liegt, ist die Festplattenaktivität vernachlässigbar. Der Netzwerkverkehr
resultiert aus der lokalen HTTP-Kommunikation mit der Ollama-API.

% ============================================================
\section{ASR-Benchmark: Parakeet TDT 0.6B}
\label{sec:eval_asr}
% ============================================================

Zur Evaluierung der Spracherkennung wurde eine Testaufnahme mit 14 Sätzen
aus dem robotikspezifischen Kontext (ROS2, igus ReBeL) verwendet.
Die Aufnahme (55\,s, 16\,kHz Mono) wurde durch das Parakeet TDT 0.6B
Modell transkribiert und die Ergebnisse satzweise gegen die Referenztranskription
ausgewertet. Als Metriken dienen die \textbf{Wortfehlerrate} (WER),
die \textbf{Zeichenfehlerrate} (CER) sowie der \textbf{Echtzeit-Faktor} (RTF).

\subsection{Wortfehlerrate pro Satz}

\begin{figure}[H]
    \centering
    \includegraphics[width=0.85\textwidth]{asr_wer_pro_satz.png}
    \caption{Wortfehlerrate (WER) pro Satz. 6 von 14 Sätzen werden fehlerfrei transkribiert.
    Die durchschnittliche WER beträgt 9,8\,\%.}
    \label{fig:eval_wer}
\end{figure}

Abbildung~\ref{fig:eval_wer} zeigt die WER pro Satz. Die durchschnittliche
Gesamt-WER beträgt \textbf{9,8\,\%}; 6 von 14 Sätzen werden fehlerfrei
transkribiert (WER\,=\,0\,\%). Die höchste Fehlerrate tritt bei Satz~8 auf
(\glqq{}Set the maximum velocity to fifty percent\grqq{}, WER\,=\,29\,\%),
da das Modell die ausgeschriebene Zahl \glqq{}fifty percent\grqq{} als
\glqq{}50\,\%\grqq{} transkribiert.

\subsection{Fehlerverteilung}

\begin{figure}[H]
    \centering
    \includegraphics[width=0.85\textwidth]{asr_fehlerverteilung.png}
    \caption{Fehlerverteilung pro Satz: Anzahl korrekter Wörter, Ersetzungen, Löschungen und Einfügungen.
    Die Mehrzahl der Wörter wird korrekt erkannt.}
    \label{fig:eval_error_dist}
\end{figure}

Die Fehleranalyse in Abbildung~\ref{fig:eval_error_dist} zeigt, dass die
dominierenden Fehlerquellen Ersetzungen sind, insbesondere bei domänentypisch fachspezifischen
Eigennamen wie \glqq{}Igus ReBeL\grqq{} (transkribiert als \glqq{}EGOS
Rebel\grqq{}) und \glqq{}MoveIt\grqq{} (transkribiert als \glqq{}MoveId\grqq{}).
Solche Fehler könnten durch ein domänenspezifisches Sprachmodell (Language Model)
oder eine Nachverarbeitungspipeline reduziert werden.

\subsection{WER vs. CER}

\begin{figure}[H]
    \centering
    \includegraphics[width=0.85\textwidth]{asr_wer_vs_cer.png}
    \caption{Vergleich von Wortfehlerrate (WER) und Zeichenfehlerrate (CER) pro Satz.
    Die CER liegt durchgehend unter der WER, was auf phonetisch korrekte Transkriptionen hindeutet.}
    \label{fig:eval_wer_cer}
\end{figure}

Abbildung~\ref{fig:eval_wer_cer} vergleicht WER und CER. Die CER liegt
mit \textbf{5,4\,\%} deutlich unter der WER (9,8\,\%), was darauf hindeutet,
dass die meisten Fehler nur einzelne Zeichen innerhalb eines Wortes betreffen
(z.\,B. \glqq{}ReBeL\grqq{} $\rightarrow$ \glqq{}Rebel\grqq{}).
Die phonetische Grundstruktur der Wörter wird weitgehend korrekt erfasst.

\subsection{Zusammenfassung der ASR-Ergebnisse}

\begin{figure}[H]
    \centering
    \includegraphics[width=0.85\textwidth]{asr_zusammenfassung.png}
    \caption{Zusammenfassung der ASR-Benchmark-Kennzahlen: Gesamt-WER 9,8\,\%, Gesamt-CER 5,4\,\%,
    Echtzeit-Faktor 0,09x (ca.\,11-fache Echtzeitverarbeitung).}
    \label{fig:eval_asr_summary}
\end{figure}

Die Zusammenfassung in Abbildung~\ref{fig:eval_asr_summary} zeigt die
wichtigsten Kennzahlen. Der Echtzeit-Faktor von \textbf{0,09x} bedeutet,
dass 55 Sekunden Audio in nur 4,9 Sekunden verarbeitet werden, was einer
ca.\,11-fachen Echtzeitverarbeitung auf der CPU entspricht.

% ============================================================
\section{Zusammenfassung der Evaluation}
% ============================================================

\begin{table}[H]
    \centering
    \small
    \begin{tabular}{l l r}
        \toprule
        Metrik & Komponente & Ergebnis \\
        \midrule
        Generierungsgeschwindigkeit & LLM (Qwen 3.5 9B) & 93,7 Tokens/s \\
        VRAM-Belegung & LLM & $\approx$\,8\,GB \\
        GPU-Leistungsaufnahme & LLM & $\approx$\,320\,W \\
        \midrule
        Gesamt-WER & ASR (Parakeet TDT) & 9,8\,\% \\
        Gesamt-CER & ASR & 5,4\,\% \\
        Echtzeit-Faktor (RTF) & ASR & 0,09x \\
        Perfekte Sätze & ASR & 6/14 (43\,\%) \\
        \bottomrule
    \end{tabular}
    \caption{Zusammenfassung der Evaluationsergebnisse.}
    \label{tab:eval_zusammenfassung}
\end{table}

Die Evaluation zeigt, dass die eingesetzte On-Premise-Lösung sowohl bei der
LLM-Inferenz als auch bei der Spracherkennung praxistaugliche Ergebnisse liefert.
Das Qwen 3.5 9B Modell erreicht eine hohe Generierungsgeschwindigkeit von über
90 Tokens/s auf der RTX 3080 Ti, während das Parakeet TDT 0.6B ASR-Modell
eine Gesamt-WER von unter 10\,\% bei 11-facher Echtzeitverarbeitung erzielt.
Die hier gewonnenen Evaluationsergebnisse werden im folgenden Kapitel um
Sicherheits- und Fehlerbetrachtungen ergänzt, die für den produktiven Einsatz
in der Robotik entscheidend sind.
